% SOURCE_AUTHORITY: sga5_fr_workpass.tex, lines 12870--12894 (LF-normalized unit).
% SOURCE_SHA256: 791F4EFFC5E02832D5D77ED03518C8156D6F07E4C8238B03545DB93D883FBB28
De hecho, el objeto de $\Delta$ definido por $\Hom^\bullet_{\Lambda[G]}(Sw_x^\Lambda,M)$ no depende, salvo isomorfismo, del recubrimiento de Galois $C'\to C$ que trivializa $F|U$. Para demostrarlo, basta comparar los elementos asociados a dos recubrimientos
\[
C'\xrightarrow{p}C,\qquad C''\xrightarrow{p'}C'\xrightarrow{p}C.
\]
En este caso, si $G'$ es el grupo de Galois del recubrimiento $C''\xrightarrow{p''}C$, entonces $G$ es un grupo cociente de $G'$. Sea
\[
\varphi:\mathbb Z_\ell[G']\longrightarrow \mathbb Z_\ell[G]
\]
el epimorfismo canónico asociado. Los módulos de Swan están entonces ligados por la relación
\[
Sw_{G,x}\simeq Sw_{G',x}\otimes_{\mathbb Z_\ell[G']}\mathbb Z_\ell[G],
\]
que resulta inmediatamente del hecho de que el carácter $a_{G,x}$ es el carácter inducido de $G'$ a $G$ ([3] VI Prop. 3).

Por otra parte, se tiene un isomorfismo de complejos de $A[G']$-módulos
\[
F_{\bar\eta''}\simeq [F_{\bar\eta'}]_{(\varphi)}.
\]
De ello resulta un isomorfismo canónico de complejos de $A$-módulos entre
\[
\Hom^\bullet_{\Lambda[G]}(Sw_{G,x}^\Lambda,F_{\bar\eta})
\quad\text{y}\quad
\Hom^\bullet_{\Lambda[G']}(Sw_{G',x}^\Lambda,F_{\bar\eta''})
\]
(fórmula de adjunción).
